\chapter{Objetivos y estado general}

\section*{Objetivo general}

Evaluar el desempeño de generalización de máquinas de aprendizaje profundo
frente al fenómeno de cambio de dominio en sistemas BCI multiusuario de
imaginación motora, mediante un protocolo de validación inter-sujeto que
permita la identificación de la arquitectura con mayor robustez.

\section*{Objetivos específicos}

\begin{enumerate}
  \item Seleccionar un conjunto de datos multisujeto de señales sEEG\footnote{En este
        documento, sEEG designa el electroencefalograma de superficie, registrado sobre el
        cuero cabelludo, y no la estereoelectroencefalografía invasiva.} asociadas
        a tareas de imaginación motora, que cuente con la representación y
        diversidad inter-sujeto necesarias para el entrenamiento y evaluación
        de las máquinas de aprendizaje profundo.
        \label{oe:1}

  \item Aplicar un flujo de preprocesamiento en las señales sEEG para la
        eliminación de ruido y artefactos, selección de canales espaciales,
        obtención de bandas de frecuencia y transformación a representaciones
        de entrada compatibles con las arquitecturas seleccionadas.
        \label{oe:2}

  \item Implementar y configurar las arquitecturas Conformers, Mamba y
        \textit{Foundation Models} para la decodificación de imaginación
        motora.
        \label{oe:3}

  \item Evaluar comparativamente el desempeño de generalización de las
        arquitecturas implementadas mediante un protocolo de validación
        inter-sujeto y un indicador de degradación por cambio de dominio.
        \label{oe:4}
\end{enumerate}

% Sin esto, el título queda al pie de una página y la tabla empieza en la
% siguiente. Un `longtable` no se puede mover, así que se reserva el sitio
% antes: si no caben 10 líneas, salta de página el conjunto.
\needspace{10\baselineskip}
\section*{Estado de avance}

La tabla~\ref{tab:estado} resume el estado de avance de cada objetivo específico.

\begin{longtable}{@{}p{0.8cm}p{2.2cm}p{10.5cm}@{}}
\caption{Estado de avance por objetivo específico.}\label{tab:estado}\\
\toprule
\textbf{OE} & \textbf{Estado} & \textbf{Descripción del avance} \\
\midrule
\endfirsthead
\toprule
\textbf{OE} & \textbf{Estado} & \textbf{Descripción del avance} \\
\midrule
\endhead

1 & Completo & Se identificaron las bases de datos públicas disponibles y se
aplicó un proceso de inclusión y exclusión, comparándolas por número de clases,
cantidad de sujetos y paradigma de imaginación motora. Se excluyeron además los
conjuntos que declaran como preentrenamiento los modelos fundacionales, para
reducir el riesgo de contaminación entre preentrenamiento y evaluación. Se seleccionaron dos
conjuntos con papeles distintos: \openbmi{} como principal y WBCIC-SHU para la
validación externa. \\
\addlinespace

2 & Completo, con salvedades & Se definieron las técnicas de preprocesamiento:
filtrado en una banda única de \hz{0.5} a \hz{40}, detección y eliminación de
artefactos, segmentación y normalización. No se hizo selección de canales con un
criterio propio: se emplea el montaje completo y, en la validación externa, la
intersección de los dos montajes. La comparación de subbandas, ventanas y aumentación se realizó en el protocolo intra-sujeto con
la semilla 2026 (sección~\ref{sec:ablaciones-f4}): en \eegnet{} y el Conformer con ocho variantes
y en LaBraM y EEGMamba solo con la ventana de \seg{0} a \seg{2}, la reflexión de canales y el
ruido en Fourier, de modo que las subbandas no se evaluaron en estas dos. Los efectos dependen de
la arquitectura y no se confirmaron con más semillas. Se especificaron las representaciones de entrada de cada arquitectura y se
verificó la lectura del conjunto de datos contra la literatura, replicando un resultado
publicado con la banda, los canales y la ventana de sus autores. \\
\addlinespace

3 & Completo, con salvedades & Está fijado el esquema de entrenamiento unificado para todas las
arquitecturas. Se implementaron y evaluaron las cuatro: la arquitectura de
referencia (\eegnet{}), el Conformer, el modelo fundacional LaBraM (variante Base) y
EEGMamba, sobre las mismas particiones y con los dos protocolos, inter-sujeto e
intra-sujeto. EEGMamba requirió adaptaciones específicas
(sección~\ref{sec:eegmamba-adapt}): con la cabeza de promediado global su desempeño
intra-sujeto quedaba apenas por encima del azar, y la cabeza que conserva la
representación por canal lo sitúa claramente por encima de él, con lo que la configuración
queda fijada para ambos protocolos. Las cuatro arquitecturas se ejecutaron con tres semillas (2026, 2027 y 2028) en los protocolos
intra-sujeto y \loso{} y en la transferencia (sección~\ref{sec:semillas-por-arq}). Las semillas no
comparten del todo el régimen de optimización, pues el recorte de norma de gradiente no estuvo
presente en todas (tabla~\ref{tab:regimen}), y esa condición se trata en las
secciones~\ref{sec:recorte-gradiente} y~\ref{sec:confusor}. \\
\addlinespace

4 & En curso & Están implementadas las métricas de desempeño y el protocolo de
validación inter-sujeto. Se propone sustituir el indicador de degradación
descrito en la literatura por una formulación que aísla el cambio de dominio
del volumen de datos de entrenamiento. Se añadió la validación sobre un segundo
conjunto de datos, que descompone la degradación en la parte atribuible a la persona
y la atribuible al protocolo de adquisición
(capítulo~\ref{cap:validacion-externa}). Con los resultados de tres semillas no se identifica
una única arquitectura de mayor robustez: ninguna obtiene el mejor valor en los tres ejes de la
definición (sección~\ref{sec:criterio-robustez}), pues el menor \dsi{} corresponde a \eegnet{} y el
menor \idp{} a EEGMamba. La diferencia de \idp{} queda dentro de la dispersión entre semillas solo
entre el Conformer y EEGMamba ($0.144$ frente a $0.137$, con una desviación entre semillas de $0.018$
en el Conformer); entre LaBraM y EEGMamba ($0.167$ frente a $0.137$) no, aunque los intervalos jerárquicos, de amplitud $0.16$ a $0.21$, no permiten
distinguir una arquitectura de otra (sección~\ref{sec:aplicacion-robustez}). Se evaluaron además el alineamiento euclidiano, la reflexión
de canales y la calibración con pocos ensayos del sujeto no visto (secciones~\ref{sec:f6f8}, \ref{sec:ea-transferencia} y~\ref{sec:fewshot}); con
los procedimientos fijados de antemano, el alineamiento no redujo la caída por protocolo y la
calibración no mejoró la exactitud. \\

\bottomrule
\end{longtable}

% =====================================================================
%  METODOLOGÍA
% =====================================================================
